% XCharter in T1, whose TFMs put the f-ligatures in slots 27 to 31 in
% another order than T1's: 29 ff, 31 fi, 30 fl, 28 ffi, and 27 ffl. Each
% line below is one of the package's fonts in T1, with every ligature, and
% the page shows the letters the PDF shows. The small capitals, the sans
% serif, and the typewriter text have no such ligatures, and guard them.
\documentclass{article}
\usepackage[T1]{fontenc}
\usepackage{XCharter}
\usepackage{amsthm}
\usepackage{hyperref}
\newtheorem{theorem}{Theorem}
\begin{document}
\section{Baffled officers fly}

Upright: off first buffer efficient fly baffle.

\textit{Italic: off first buffer efficient fly baffle.}

\textbf{Bold: off first buffer efficient fly baffle.}

\textbf{\textit{Bolditalic: off first buffer efficient fly baffle.}}

\textsl{Slanted: off first buffer efficient fly baffle.}

\textbf{\textsl{Boldslanted: off first buffer efficient fly baffle.}}

\textsc{Smallcaps: off first buffer efficient fly baffle.}

Marked text.\footnote{Note: off first buffer efficient fly baffle.}

\begin{theorem}
Theorem: off first buffer efficient fly baffle.
\end{theorem}

{\fontfamily{XCharter-LF}\selectfont Lining: off first buffer efficient fly baffle.}

{\fontfamily{XCharter-OsF}\selectfont Oldstyle: off first buffer efficient fly baffle.}

{\fontfamily{XCharter-TOsF}\selectfont Tabularoldstyle: off first buffer efficient fly baffle.}

{\fontfamily{XCharter-TLF}\selectfont Tabularlining: off first buffer efficient fly baffle.}

\textsf{Sans: off first buffer efficient fly baffle.}

\texttt{Mono: off first buffer efficient fly baffle.}
\end{document}
